\documentclass[12 pt]{report}

\usepackage[utf8]{inputenc}
\usepackage[T1]{fontenc}

\usepackage{amssymb,amsmath}
\usepackage{textcomp}
\usepackage{mathtools}
\usepackage{enumitem}
\usepackage{bm}	% In order to have bold italic in math mode
\usepackage{todonotes}	% In order to have clean todo's
\usepackage[labelfont=bf]{caption}
\usepackage{subcaption}
\usepackage{graphicx}
\usepackage{epstopdf}
\usepackage[colorlinks=true,linkcolor=red]{hyperref}
\usepackage{cleveref}


\begin{document}

\begin{titlepage}
    \vspace{-1cm}
    
    \begin{minipage}{0.6\textwidth}
        \begin{flushleft}
            \centering
            \textsc{University of Liège} \\
            \textsc{Faculty of Applied Sciences} \\
            \textsc{Academic Year 2017 - 2018}
        \end{flushleft}
    \end{minipage}

    \vspace{-2.6cm}
    \hspace{0.5\textwidth}
    
    \begin{minipage}[t]{\textwidth}
        \begin{flushright}
            \includegraphics[scale = 0.6]{uLIEGE_Faculte_SciencesAppliquees_Logo_CMJN_pos}
        \end{flushright}
    \end{minipage}

    \vspace{1cm}


    \begin{minipage}{\textwidth}
        \hspace{-0.7cm}\noindent\makebox[\linewidth]{\rule{\textwidth}{2pt}} 
    \end{minipage}

    \vspace{1cm}

    \begin{minipage}{\textwidth}
        \begin{center}
           \hspace{-0.8cm}\Huge{\textsc{\textbf{MATH0471-Multiphysics Integrated Computational Project}}}\\
            \vspace{0.3cm}
            \hspace{-0.8cm}\textsc{C. GEUZAINE \\ R. Boman} \\
        \end{center}
    \end{minipage}

    \vspace{1cm}

    \begin{minipage}{\textwidth}
        \hspace{-0.7cm}\noindent\makebox[\linewidth]{\rule{\textwidth}{0.2pt}}
    \end{minipage}

    \vspace{1cm}

    \begin{minipage}{\textwidth}
        \begin{center}
            \hspace{-1cm}\Huge{\textbf{Intermediate report 1}} \\ 
            \vspace{0.5cm}
        \end{center}
    \end{minipage}

    \vspace{1cm}

    \begin{minipage}{\textwidth}
        \hspace{-0.7cm}\noindent\makebox[\linewidth]{\rule{\textwidth}{0.2pt}}
    \end{minipage}

    \vspace{1cm}

    \begin{minipage}{\textwidth}
        \begin{center}
            \hspace{-1cm}\large{\textsc{DELCOURT Arnaud}} \\
            \hspace{-1cm}\large{\textsc{DESSERS Christophe}} \\
            \hspace{-1cm}\large{\textsc{HEUCHAMPS Alexandre}} \\
            \hspace{-1cm}\large{\textsc{TOMASETTI Romin}} \\
            \vspace{1cm}
            \hspace{-1cm}\Large{\textsc{\today}} \\
        \end{center}
    \end{minipage}
    
    \vspace{1cm}
    
    \begin{minipage}{\textwidth}
        \hspace{-0.7cm}\noindent\makebox[\linewidth]{\rule{\textwidth}{2pt}}
    \end{minipage}
\end{titlepage}

\tableofcontents
\newpage
\listoffigures
\newpage
\listoftables
\newpage

\section{Questions to ask on February, 20}

\begin{itemize}
	%
	\item Do we have to take into account the frequency-dependency of the electromagnetic properties of the materials? (see \url{https://em.geosci.xyz/content/physical_properties/magnetic_permeability/magnetic_permeability_frequency_dependent.html})
	%
	\item Regarding the already-written code:
		\begin{itemize}
			\item Comments on the created classes?
			\item Is speed affected? Or we don't care?
			\item Especially, ask if Array\_3D\_Template<Node> nodes is a good way to start !
		\end{itemize}
	%
	\item Regarding testing, what are the simplest tests? Could we compare our results with other simulations? Do we have to compare with something that can be hand-solved?
	%
	\item Check our reasoning on the nodes. They all have:
		\begin{itemize}
			\item 3 coordinates (3 doubles) - How are we going to mesh ??? (we think we are obliged to keep rectangular parallelepiped because of the numerical scheme)
			%
			\item 6 fields (3 for $\vec{E}$ and 3 for $\vec{H}$ so 6 doubles)
			%
			\item 1 temperature (1 double)
			%
			\item 1 unsigned char for the ID of the material
			%
			\item $\Rightarrow$ That would make 10 doubles and one unsigned char or 10*4+1 bytes per node.
		\end{itemize}
\end{itemize}

\section{Database for human tissues}

\url{https://www.itis.ethz.ch/virtual-population/tissue-properties/downloads/database-v3-1/}

\end{document}
